\documentclass[../main.tex]{subfiles}

\begin{document}
	\chapter{Modelli per lo studio della viscosità}
	% Tutto quello che serve per capire i capitoli dopo
	
	\section{Liquidi sottoraffreddati e viscosità}
	
	Un liquido si definisce "sottoraffreddato" nel momento in cui viene portato al di sotto della sua temperatura di fusione $T_m$ evitando il processo di cristallizzazione.
	Man mano che si diminuisce la temperatura, mantenendo il sistema in questo stato metastabile, la viscosità del liquido aumenta di diversi ordini di grandezza fino ad arrivare al punto critico in cui il tempo di rilassamento supera il tempo di osservazione sperimentale, per cui la sua evoluzione appare congelata (Fig. \ref {glass_transition_plot}).
	Questa è quella che viene definita transizione vetrosa, a seguito della quale il liquido si trasforma in vetro. Non si tratta di una vera e propria trasformazione di fase termodinamica, in quanto microscopicamente si ha ancora a che fare con un liquido sottoraffreddato dall'elevata viscosità, ma macroscopicamente il materiale manifesta le proprietà di un solido.
	La temperatura di transizione vetrosa $T_g$ dipende dalla velocità di raffreddamento e dal tempo di osservazione sperimentale preso come riferimento. In seguito verrà definita, seguendo la convenzione standard, come la temperatura alla quale la viscosità $\eta (T_g) = 10^{12}$ Pa s \cite{cavagnaSupercooledLiquidsPedestrians2009}.
	
	\begin{figure}[h!]
		\centering
		\includegraphics[height=8cm]{entropy_temperature_plot.png}
		\caption{Rappresentazione schematica dell'entropia di un liquido in funzione della temperatura \cite{cavagnaSupercooledLiquidsPedestrians2009}.}
		\label{glass_transition_plot}
	\end{figure}
	
	L'andamento della viscosità in funzione della temperatura varia al variare del liquido sotto esame come mostrato in Fig. \ref{angell_plot}. Più la curva si discosta da una retta e più l'indice di fragilità $m$ aumenta. Quest'ultimo è definito come la pendenza di $log(\eta)$ in funzione di $T/T_g$ nel punto $T/T_g = 1$ \cite{cavagnaSupercooledLiquidsPedestrians2009}. Liquidi con bassi valori di $m$ (come il vetro delle finestre $SiO_2$) vengono detti forti, mentre quelli aventi alti valori di $m$ fragili.
	
	\begin{figure}[h!]
		\centering
		\includegraphics[height=10cm]{angell_plot.png}
		\caption{Plot di Angell - Andamento del logaritmo della viscosità in funzione dell'inverso della temperatura \cite{angellFormationGlassesLiquids1995}.}
		\label{angell_plot}
	\end{figure}
	
	L'intercetta del plot di Angell indica invece il valore teorico assegnato alla viscosità nel limite in cui la temperatura tende all'infinito. Di seguito verrà indicata con $\eta_\infty = \lim_{T->\infty} \eta(T)$. Tale valore, sebbene fisicamente irrilevante, viene introdotto in quanto spesso usato come parametro dalle forme funzionali empiriche che cercano di interpolare l'andamento della viscosità dei liquidi sottoraffreddati (di seguito chiamati anche "liquidi formatori di vetri") in funzione della temperatura.
	Tra queste la più nota è la legge di VFT (Vogel-Fulcher-Tamman) \cite{fulcherAnalysisRecentMeasurements1925}:
	\begin{equation}
		\log_{10} \eta(T) = \log_{10} \eta_\infty + \frac{(12 - \log_{10} \eta_\infty)^2}{m (T/T_g - 1) + (12 - \log_{10} \eta_\infty)}
	\end{equation}
	Un altra ancora è l'equazione MYEGA (Mauro–Yue–Ellison–Gupta–Allan), che a basse temperature risulta più accurata della prima \cite{mauroViscosityGlassformingLiquids2009}:
	\begin{equation}
		\begin{aligned}
			&\log_{10} \eta(T) \\
			&= \log_{10} \eta_\infty + (12 - \log_{10} \eta_\infty) \frac{T_g}{T} \exp \left[ \left( \frac{m}{12 - \log_{10} \eta_\infty} - 1 \right) \left( \frac{T_g}{T} - 1 \right) \right]
		\end{aligned}
	\end{equation}


	\section{Modelli data-driven}
		Perché ci interessa tutto questo?? Discorso introduttivo di Cassar...
		Features in input- output e perché vogliamo fare questo fit
		svantaggi e vantaggi, approcci
	
	\section{Perché un modello lineare?}
	%TODO: da riscrivere meglio (è la prima sezione che ho buttato giù)
	
	Lo standard attuale per la realizzazione di modelli data driven nello studio dei liquidi sottoraffreddati si basa sull'utilizzo di neural networks per generare predizioni delle loro proprietà partendo della loro composizione chimica. Tale procedura consente di ottenere risultati molto accurati grazie alla elevata capacità delle reti neurali di individuare pattern complessi tra features (input) e targets (output) con un set di dati su cui allenare la rete sufficientemente ampio.
	Tuttavia, il problema di questo approccio risiede nella bassa interpretabilità dei risultati ottenuti. Proprio per questo motivo questi modelli vengono chiamati "black-box". Si possono ricavare alcune informazioni sul loro funzionamento interno tramite analisi computazionalmente impegnative come la SHAP (SHapley Additive exPlanations), tuttavia tale procedura è in grado di fornire solo informazioni indicative e non sempre accurate sull'interpretabilità delle predizioni.
	Un metodo interessante per migliorare l'interpretabilità di questo tipo di modelli è quello adottato da Cassar \autocite{cassarViscNetNeuralNetwork2021} nella realizzazione di ViscNet. Tale modello mira a studiare con un approccio "gray-box" la dipendenza della viscosità dalla temperatura di "vetri basati su ossidi". Al posto di far predire al modello direttamente la viscosità, la rete neurale predice i parametri dell'equazione MYEGA: l'indice di fragilità del liquido, la temperatura di transizione vetrosa e la viscosità asintotica. Grazie a questi parametri, risulta poi possibile calcolare la viscosità a qualsiasi temperatura di interesse sfruttando il vantaggio delle buone capacità di estrapolazione intrinseche dell'equazione.
	Sebbene questo fornisca informazioni in più rispetto ad un approccio puramente black-box, il problema intrinseco sull'interpretabilità della NN permane.
	Per tale ragione, nello svolgimento di questo lavoro di tesi, si è scelto di replicare l'analisi fatta da Cassar utilizzando un modello lineare. Quest'ultimo, infatti, in linea di principio offre totale trasparenza sui meccanismi interni che collegano features e targets. Naturalmente questo è inevitabilmente correlato a una perdita di performance siccome le dipendenze fisiche tra i dati spesso non sono lineari. Tuttavia non è sempre il caso, per cui vale comunque la pena tentare di esplorare in via preliminare la situazione con un modello più semplice prima di adottare una rete neurale anche nell'ottica di effettuare predizioni. In questo modo, tramite lo studio dei pesi interni del modello lineare, risulta possibile ricavare informazioni preziose riguardanti la fisica che sta dietro ai risultati ottenuti.
	%TODO:	spiegare meglio la Filosofia dello stuido?
	
	\subsection{Confronto tra regressione lineare e reti neurali}
	%! Serve davvero??
	%	più facili da intepretari, sono più seplici, in ambito fisico si prestano ad un intepretazione tra features in input e output
	
	
	\subsection{Il problema della multicollinerità}
	
	
	\subsection{Ridge Regression}
	In particolare è stata utilizzata una ridge regression durante il corso dell'analisi per poter limitare gli effetti della mulitcollinearità tra le features che ha come conseguenza la perdita del senso fisico dei pesi del modello.
	Spiegare perchè...
	
\end{document}